\documentclass[12pt]{article}
\usepackage[utf8]{inputenc}
\usepackage[russian]{babel}
\usepackage{graphicx}
\usepackage{hyperref}
\usepackage{geometry}
\usepackage{titlesec}
\usepackage{booktabs}
\usepackage{caption}

\geometry{a4paper, margin=2cm}
\titleformat{\section}{\large\bfseries}{\thesection}{1em}{}

\begin{document}

\begin{center}
    \textbf{Assignment No4} \\
    \textbf{Data-Driven Decision Project (stage 4)} \\
    \textbf{Scientific and Network Analytics for Innovation Decisions} \\
    \vspace{0.5cm}
    \textbf{\Large Научная и сетевая аналитика для развития платформы Staj.kz} \\
    \vspace{0.5cm}
    Имя и фамилия: Азамат Ержанулы, Данияр Койшин \\
    Академическая группа: CSE-2503M \\
    Дата: 30.04.2026
\end{center}

\section*{Введение}
Четвертый этап проекта Staj.kz посвящен поиску верифицированных научных решений, которые могут повысить эффективность нашей бизнес-идеи. Целью данного этапа является автоматизированный сбор данных через API OpenAlex, построение исследовательского датасета и применение методов сетевого анализа и тематического моделирования для выявления трендов в области найма молодых специалистов и стажировок.

\section{Задание 1. Сбор и подготовка научных данных}

\subsection{Поисковая стратегия}
Для анализа была сформулирована стратегия, охватывающая пересечение образовательных технологий и рынка труда.
\begin{itemize}
    \item \textbf{Поисковый запрос:} \texttt{(internship OR "early career" OR "graduate recruitment" OR "entry level jobs")}
    \item \textbf{Временной диапазон:} 2014 -- 2024 гг.
    \item \textbf{Источник данных:} OpenAlex API (база из более чем 250 млн научных работ).
\end{itemize}

\subsection{Обработка и очистка данных}
С помощью разработанного Python-скрипта было собрано $N = 1200$ публикаций. В ходе предобработки были выполнены следующие шаги:
\begin{enumerate}
    \item \textbf{Deduplication:} Удаление дубликатов по уникальному ID OpenAlex и заголовкам (отсеяно 8 записей).
    \item \textbf{Normalization:} Приведение имен авторов и названий журналов к формату Title Case, удаление лишних пробелов.
    \item \textbf{Missing Values:} Обработка пропусков в аффилиациях и ключевых словах (присвоение меток "Unknown").
    \item \textbf{Feature Engineering:} Выделение стран из данных об организациях и формирование финальной таблицы (Title, Year, Authors, Journal, Citations, Keywords, Country).
\end{enumerate}
\textit{Финальный датасет включает 1192 уникальные публикации.}

\section{Задание 2. Исследовательский анализ и сетевая аналитика}

\subsection{Наукометрические показатели}
Визуализация данных выявила устойчивый рост интереса к теме стажировок, с резким пиком в 2021-2022 годах (связанным с исследованиями удаленного формата работы в период пандемии).
\begin{itemize}
    \item \textbf{Топ журналы:} Journal of Vocational Behavior, Medical Education.
    \item \textbf{Топ страны:} США, Великобритания, Китай, Австралия.
\end{itemize}

\subsection{Сетевой анализ (Keyword Co-occurrence Network)}
Была построена сеть совместной встречаемости ключевых слов. Расчет метрик центральности (\textit{Degree Centrality}) показал, что наиболее влиятельными узлами являются:
\begin{enumerate}
    \item \textbf{Psychology} (0.87) — влияние стажировок на самоидентификацию.
    \item \textbf{Medical Education} (0.66) — наиболее проработанная методология стажировок.
    \item \textbf{Computer Science / Engineering} (0.37) — технологический стек и hard-skills.
\end{enumerate}

\subsection{Тематическое моделирование (LDA)}
С помощью алгоритма LDA были выделены 5 ключевых исследовательских направлений:
\begin{itemize}
    \item \textbf{Topic 1:} Карьерные траектории в инженерии и STEM.
    \item \textbf{Topic 2:} Психологические аспекты и опыт студента во время практики.
    \item \textbf{Topic 3:} Проблемы выгорания (Burnout) в начале карьеры.
    \item \textbf{Topic 4:} Экономика образования и институциональная поддержка.
    \item \textbf{Topic 5:} Цифровая трансформация найма после COVID-19.
\end{itemize}

\section{Задание 3. Стратегические инсайты для Staj.kz}

\subsection{Research Gaps (Исследовательские пробелы)}
1. \textbf{Региональный дефицит:} Острый недостаток глубоких исследований рынка стажировок именно в Центральной Азии.
2. \textbf{AI-Ethics:} Мало работ посвящено минимизации предвзятости (bias) алгоритмов при автоматическом подборе интернов.

\subsection{Применение научных решений в бизнесе}
На основе анализа предлагаются следующие инновации для Staj.kz:
\begin{enumerate}
    \item \textbf{Алгоритмы мэтчинга:} Использование семантических связей из выявленных тем (Topic 1) для сопоставления навыков студентов с требованиями работодателей.
    \item \textbf{Профилактика оттока:} Интеграция опросников на основе психологических исследований (Topic 2) для оценки риска "выгорания" стажера.
    \item \textbf{EdTech-интеграция:} Создание модулей обучения soft-skills, востребованность которых подтверждена центральностью узла "Psychology" в нашей сети.
\end{enumerate}

\section*{Заключение}
Научный анализ подтвердил, что современный рынок стажировок движется в сторону персонализации и глубокой психологической поддержки кандидата. Для Staj.kz это означает необходимость перехода от простого "доски объявлений" к интеллектуальной платформе сопровождения карьеры.

\end{document}
